% META: {"id": "TSPE-DERIVATION-CONVEXITE-EV-A", "chapitre": "TSPE-DERIVATION-CONVEXITE", "type_objet": "evaluation", "version": "A", "duree_min": 55, "bareme_total": 20, "capacites_codes": ["C1", "C2", "C3", "C4", "C5", "C6"], "competences": ["chercher", "modeliser", "representer", "calculer", "raisonner", "communiquer"], "mode_creation": "ex_nihilo", "sources_inspiration": [], "status": "approved"}
% BEGIN-VERIFY
% from sympy import *
% x = symbols('x')
%
% # ---- Exercice 1 (5 pts) : Derivee de composee (C1) ----
% f1 = exp(2*x**2 + 1)
% assert diff(f1, x) == 4*x*exp(2*x**2 + 1)
% f2 = ln(x**2 + 3)
% assert simplify(diff(f2, x)) == 2*x/(x**2 + 3)
%
% # ---- Exercice 2 (5 pts) : Etude complete (C2) ----
% g = x**2 * exp(-x)
% gp = diff(g, x)
% assert simplify(gp) == (2 - x)*x*exp(-x)
% assert solve(gp, x) == [0, 2]
% assert g.subs(x, 2) == 4*exp(-2)
% assert limit(g, x, oo) == 0
% assert limit(g, x, -oo) == oo
%
% # ---- Exercice 3 (5 pts) : Convexite et inegalites (C3, C6) ----
% h = exp(x) - 1 - x
% assert diff(h, x, 2) == exp(x)  # > 0, convexe
% assert h.subs(x, 0) == 0
% assert diff(h, x).subs(x, 0) == 0
% # h convexe + min en 0 => h >= 0 => e^x >= 1+x
%
% # ---- Exercice 4 (5 pts) : Esquisse et inflexion (C4, C5) ----
% k = x**3 - 6*x**2 + 9*x + 1
% kp = diff(k, x)
% kpp = diff(k, x, 2)
% assert expand(kp) == 3*x**2 - 12*x + 9
% assert expand(kpp) == 6*x - 12
% assert solve(kpp, x) == [2]
% assert k.subs(x, 2) == 3
% assert kp.subs(x, 2) == -3
% # tangente en 2: y = -3x + 9, k traverse sa tangente (inflexion)
% END-VERIFY

\textbf{Évaluation A -- TSPE-DERIVATION-CONVEXITE}

\begin{center}
\textit{Duree : 55 min}
\baremeIndicatif{20 points}
\end{center}

\textbf{Exercice 1.} \textit{(5 points)} -- \textbf{Dérivée de fonctions composeees (C1)}

Calculer les dérivées des fonctions suivantes en precisant le domaine de dérivation :
\begin{enumerate}
  \item $f_1(x) = \mathrm{e}^{2x^2 + 1}$ sur $\mathbb{R}$.
  \item $f_2(x) = \ln(x^2 + 3)$ sur $\mathbb{R}$.
\end{enumerate}

\bigskip

\textbf{Exercice 2.} \textit{(5 points)} -- \textbf{Étude complete (C2)}

Soit $g(x) = x^2\,\mathrm{e}^{-x}$ définie sur $\mathbb{R}$.
\begin{enumerate}
  \item Calculer $g'(x)$ et étudier son signe.
  \item Déterminer les limites de $g$ en $+\infty$ et $-\infty$.
  \item Dresser le tableau de variations et preciser l'extremum.
\end{enumerate}

\bigskip

\textbf{Exercice 3.} \textit{(5 points)} -- \textbf{Convexité et inégalités (C3, C6)}

\begin{enumerate}
  \item Soit $h(x) = \mathrm{e}^x - 1 - x$. Calculer $h''(x)$ et en déduire la convexité de $h$.
  \item Déterminer le minimum de $h$ et en déduire l'inégalité $\mathrm{e}^x \geq 1 + x$.
  \item En utilisant la convexité de $\mathrm{e}^x$, déterminer la position de la courbe par rapport a sa tangente en $0$.
\end{enumerate}

\bigskip

\textbf{Exercice 4.} \textit{(5 points)} -- \textbf{Esquisse et points d'inflexion (C4, C5)}

Soit $k(x) = x^3 - 6x^2 + 9x + 1$.
\begin{enumerate}
  \item Calculer $k'(x)$ et $k''(x)$.
  \item Déterminer les intervalles de convexité et le point d'inflexion.
  \item Déterminer l'équation de la tangente au point d'abscisse $2$ et preciser si la courbe la traverse.
\end{enumerate}
